\chapter{PETSCII-codes}
\label{appendix:petsciicodes}

\section{PETSCII-codes}

PETSCII is de tekencodering die Commodore-8-bitcomputers gebruiken voor gegevens. De codering omvat afdrukbare tekens zoals letters, cijfers, leestekens en grafische symbolen; Controlcodes om de weergaveparameters te wijzigen en de cursor van de schermeditor te bedienen; en voorstellingen van toetsaanslagen die via het toetsenbord kunnen worden ingevoerd. Een programma kan PETSCII-codes gebruiken voor uitvoer naar het scherm, voor invoer via het toetsenbord of voor het lezen en schrijven van tekengegevens in bestanden of via seriële communicatie.

In BASIC aanvaardt de functie \textbf{CHR\$()} een PETSCII-code en geeft die code als string terug. Zo drukt \texttt{PRINT CHR\$(65)} de letter \screentext{A} af. De functie \textbf{ASC()} aanvaardt een string van één teken en geeft het getal van de PETSCII-code terug. De commando's \textbf{GET} en \textbf{GETKEY} lezen een toetsaanslag van het toetsenbord en geven de bijbehorende PETSCII-code terug.

Machinetaalprogramma's kunnen de GETIN-routine van de KERNAL of het PETSCIIKEY-hardwareregister gebruiken om toetsenbordinvoer als PETSCII-codes te lezen.

\index{Toetsenbord!PETSCII-codes en CHR\$}
\begin{adjustwidth}{}{-2cm}
\begin{multicols}{4}
\begin{description}[align=left,labelwidth=0.2cm]
    \item [0]    \item [1]   \small{Alternatief palet}    \item [2]   \small{Onderstrepen aan}    \item [3]    \item [4]   \small{Standaardpalet}    \item [5]   \small{Wit}    \item [6]    \item [7]   \small{Bel}    \item [8]    \item [9]   \specialkey{TAB}    \item [10]  \small{Nieuwe regel}
%   \item [11]  \footnotesize{DISABLE \specialkey{SHIFT}\megasymbolkey}
    \item [11]  Schakelt \\ \specialkey{SHIFT}\megasymbolkey uit
%   \item [12]  \footnotesize{ENABLE \specialkey{SHIFT}\megasymbolkey}
    \item [12]  Schakelt \\ \specialkey{SHIFT}\megasymbolkey in    \item [13]  \specialkey{RETURN}    \item [14]  \small{Kleine letters}    \item [15]  \small{Knipperen aan}    \item [16]  F9    \item [17]  \megakey{$\downarrow$}    \item [18]  \specialkey{RVS ON}    \item [19]  \specialkey{CLR\\HOME}    \item [20]  \specialkey{INST\\DEL}    \item [21]  F10 / Vorig woord    \item [22]  F11    \item [23]  F12 / Volgend woord    \item [24]  Tab instellen/wissen    \item [25]  F13    \item [26]  F14 / Vorig tabpunt    \item [27]  \small{ESCAPE}    \item [28]  \small{Rood}    \item [29]  \megakey{$\rightarrow$}    \item [30]  \small{Groen}    \item [31]  \small{Blauw}    \item [32]  \megakey{SPACE}    \item [33]  !    \item [34]  "    \item [35]  \#    \item [36]  \$    \item [37]  \%    \item [38]  \&    \item [39]  '    \item [40]  (    \item [41]  )    \item [42]  *    \item [43]  +    \item [44]  ,    \item [45]  -    \item [46]  .    \item [47]  /    \item [48]  0    \item [49]  1    \item [50]  2    \item [51]  3    \item [52]  4    \item [53]  5    \item [54]  6    \item [55]  7    \item [56]  8    \item [57]  9    \item [58]  :    \item [59]  ;    \item [60]  <    \item [61]  =    \item [62]  >    \item [63]  ?    \item [64]  @    \item [65]  A    \item [66]  B    \item [67]  C    \item [68]  D    \item [69]  E    \item [70]  F    \item [71]  G    \item [72]  H    \item [73]  I    \item [74]  J    \item [75]  K    \item [76]  L    \item [77]  M    \item [78]  N    \item [79]  O    \item [80]  P    \item [81]  Q    \item [82]  R    \item [83]  S    \item [84]  T    \item [85]  U    \item [86]  V    \item [87]  W    \item [88]  X    \item [89]  Y    \item [90]  Z    \item [91]  [    \item [92]  \pounds    \item [93]  ]    \item [94]  \megakeywhite{$\uparrow$}    \item [95]  \megakeywhite{$\leftarrow$}
\end{description}
\end{multicols}
\end{adjustwidth}

De codes 96 tot en met 127 drukken tekens af die vergelijkbaar zijn met die van de codes 192 tot en met 223. Zij worden niet als toetsenbordinvoer teruggegeven.

\begin{adjustwidth}{}{-2cm}
\begin{multicols}{4}
\begin{description}[align=left,labelwidth=0.2cm]
    \item [128]    \item [129] \small{Oranje}    \item [130] \small{Onderstrepen uit}    \item [131] \specialkey{SHIFT}\specialkey{RUN\\STOP}    \item [132] HELP    \item [133] F1    \item [134] F3    \item [135] F5    \item [136] F7    \item [137] F2    \item [138] F4    \item [139] F6    \item [140] F8    \item [141] \specialkey{SHIFT}\specialkey{RETURN}    \item [142] \small{Hoofdletters}    \item [143] \small{Knipperen uit}    \item [144] \small{Zwart}    \item [145] \megakey{$\uparrow$}    \item [146] \specialkey{RVS OFF}    \item [147] \specialkey{SHIFT}\specialkey{CLR\\HOME}    \item [148] \specialkey{SHIFT}\specialkey{INST\\DEL}    \item [149] \small{Bruin}    \item [150] \small{Lichtrood (roze)}    \item [151] \small{Donkergrijs}    \item [152] \small{Middengrijs}    \item [153] \small{Lichtgroen}    \item [154] \small{Lichtblauw}    \item [155] \small{Lichtgrijs}    \item [156] \small{Paars}    \item [157] \megakey{$\leftarrow$}    \item [158] \small{Geel}    \item [159] \small{Cyaan}    \item [160] \megakey{SPACE}    \item [161] \graphicsymbol{k}    \item [162] \graphicsymbol{i}    \item [163] \graphicsymbol{t}    \item [164] \graphicsymbol{[}    \item [165] \graphicsymbol{g}    \item [166] \graphicsymbol{<}    \item [167] \graphicsymbol{m}    \item [168] \graphicsymbol{/}    \item [169] \graphicsymbol{?}    \item [170] \graphicsymbol{m}    \item [171] \graphicsymbol{q}    \item [172] \graphicsymbol{d}    \item [173] \graphicsymbol{z}    \item [174] \graphicsymbol{s}    \item [175] \graphicsymbol{p}    \item [176] \graphicsymbol{a}    \item [177] \graphicsymbol{e}    \item [178] \graphicsymbol{r}    \item [179] \graphicsymbol{w}    \item [180] \graphicsymbol{h}    \item [181] \graphicsymbol{j}    \item [182] \graphicsymbol{l}    \item [183] \graphicsymbol{y}    \item [184] \graphicsymbol{u}    \item [185] \graphicsymbol{o}    \item [186] \graphicsymbol{\{}    \item [187] \graphicsymbol{f}    \item [188] \graphicsymbol{c}    \item [189] \graphicsymbol{x}    \item [190] \graphicsymbol{v}    \item [191] \graphicsymbol{b}    \item [192]  \graphicsymbol{C}    \item [193]  \graphicsymbol{A}    \item [194]  \graphicsymbol{B}    \item [195]  \graphicsymbol{C}    \item [196] \graphicsymbol{D}    \item [197] \graphicsymbol{E}    \item [198] \graphicsymbol{F}    \item [199] \graphicsymbol{G}    \item [200] \graphicsymbol{H}    \item [201] \graphicsymbol{I}    \item [202] \graphicsymbol{J}    \item [203] \graphicsymbol{K}    \item [204] \graphicsymbol{L}    \item [205] \graphicsymbol{M}    \item [206] \graphicsymbol{N}    \item [207] \graphicsymbol{O}    \item [208] \graphicsymbol{P}    \item [209] \graphicsymbol{Q}    \item [210] \graphicsymbol{R}    \item [211] \graphicsymbol{S}    \item [212] \graphicsymbol{T}    \item [213] \graphicsymbol{U}    \item [214] \graphicsymbol{V}    \item [215] \graphicsymbol{W}    \item [216] \graphicsymbol{X}    \item [217] \graphicsymbol{Y}    \item [218] \graphicsymbol{Z}    \item [219] \graphicsymbol{+}    \item [220] \graphicsymbol{-}    \item [221] \graphicsymbol{B}    \item [222] \graphicsymbol{\textbackslash}    \item [223] \graphicsymbol{]}

\end{description}
\end{multicols}
\end{adjustwidth}

De codes 224 tot en met 254 drukken tekens af die vergelijkbaar zijn met die van de codes 160 tot en met 190. Zij worden niet als toetsenbordinvoer teruggegeven.

De codes 126, 222 en 255 drukken alle drie het pi-symbool af (\graphicsymbol{\textbackslash}). Wanneer u dit teken typt, is de toetsenbordinvoercode 222.

Wanneer u de kleineletters-/hoofdlettermodus gebruikt (door \megasymbolkey + \specialkey{SHIFT} in te drukken), let dan op het volgende:
\begin{itemize}
  \item De hoofdletters in het gebied 65-90 van bovenstaande tabel worden vervangen door kleine letters.  \item De grafische tekens in het gebied 97-122 van bovenstaande tabel worden vervangen door hoofdletters.  \item De kleine letters van PETSCII (65-90) en de hoofdletters (97-122) vallen in het ASCII-gebied voor respectievelijk hoofdletters (65-90) en kleine letters (97-122).  \item Verscheidene symbolen hebben alternatieve symbolen in het lettertype voor kleine en hoofdletters, met name de PETSCII-codes 126, 127, 168 en 183.
\end{itemize}
